\paragraph{Measurement error.} Seven errors produced plausible but spurious
results during this study before they were caught, among them a host transfer
inside the timed region, sequential timing on a shared device, a masked
implementation benchmarked as sparse, and a hand-written sparse kernel compared
with a vendor dense kernel tuned for another regime. Appendix~\ref{sec:pitfalls}
lists each with its symptom and the practice that corrected it.

\paragraph{Limitations.} Every cost result concerns centralised inference on a
single device; the paper makes no claim about bandwidth-constrained distributed
deployment, where reducing transmitted messages is itself the objective and
top-$k$ selection may well be the right design. The analysis is single-round
and single-head. The operation-count ratio is invariant to rounds and heads, but
the timing result is not guaranteed to be: $H$ heads of width $d/H$ change the
per-head width, which is the axis along which our timing result changes sign.

The measured envelope is one GPU architecture (NVIDIA Blackwell, sm\_120), one
x86\_64 CPU, PyTorch 2.11 with CUDA 12.8 and Triton 3.6, head widths $16$ to
$512$, $N \le 512$, batch $32$ to $1024$, and fp32, TF32 and bf16. We could not replicate on a second
GPU generation, because no other was available to us. Both cards were shared
with other tenants throughout, and clocks could not be locked; every claim is
therefore a ratio between implementations timed in the same interleaved sweep,
and no claim is read off an agent count below $192$, where a fixed per-call cost
is a large share of the measurement.

Our kernels are not optimal. The sparse kernels hold the whole score row of a
query on chip, which caps them at $N \le 1024$ and admits only three tilings
above $192$ agents against $28$ to $106$ for the dense kernel; a streaming exact
top-$k$ that keeps a running candidate set would not have that constraint, and
we did not write one. The sparse efficiency figures are therefore lower bounds,
and the negative results at wide heads are weakly supported. The fused
crossover at one retained peer is provisional in the other direction, and
Section~\ref{sec:precision} shows how: a dense kernel that uses the tensor cores
removes it from $384$ agents upward. The
pre-score selectors attend to a different peer set than top-$k$ does, so their
comparison is of cost and not of quality; their selectors are composed from
library operators, so the selector costs we report are upper bounds; and their
envelope is one head width, one batch and fp32. Finally, the learning study
(Appendix~\ref{sec:training}) has five seeds per variant; it shows that the
communication architecture improves return and cannot resolve the effect of any
sparsification choice.
